% !Mode:: "TeX:UTF-8"

\chapter{FSDP通信瓶颈分析与实验验证}

\section{FSDP分布式训练框架概述}

PyTorch全分片数据并行（Fully Sharded Data Parallel, FSDP）通过切分（shard）稠密参数来扩展到单卡放不下的大模型。具体而言，FSDP把一个模型实例分解为多个更小的单元（FSDP unit）并分别处理。在前向与反向计算过程中，FSDP一次只收集一个单元的完整参数与梯度，其余单元的参数与梯度始终保持分片状态，且整个训练循环中优化器状态也保持分片\cite{zhao2023fsdp}。

FSDP的完整训练流程包含以下关键阶段：

\textbf{前向传播阶段：}首先从本地加载模型分片到GPU显存，然后各rank通过All-Gather通信收集完整参数，构建完整的FSDP单元并执行前向计算。计算完成后立即释放除本rank分片以外的参数副本，仅保留当前激活单元的显存占用。

\textbf{反向传播阶段：}再次通过All-Gather恢复完整参数，执行局部反向传播。完成后立即启动Reduce-Scatter，对梯度进行规约并重新切分，最终仅保留本rank的梯度分片，然后释放完整参数副本。

\textbf{优化器更新阶段：}在分片空间内执行本地优化器更新。

这一"分片-聚合-释放"的工作模式实现了显存占用与"分片后的模型大小$+$当前激活单元大小"近似线性相关，显著降低了单设备显存压力，同时也暴露出前向与反向阶段的All-Gather通信以及反向阶段的Reduce-Scatter通信所带来的时间开销。其中，All-Gather主要发生在参数获取阶段，可以通过预取（Prefetching）机制与后续计算进行部分重叠；而Reduce-Scatter位于反向传播的关键路径上，通常在梯度计算完成后立即启动，属于完全阻塞的通信操作，其对整体训练性能的影响更为显著。

\section{实验设置}

为确保性能数据的可比性与实验结论的统计可靠性，本研究采用控制变量法设计实验。通过固定模型架构、训练数据集、超参数配置及硬件环境等所有非实验变量，仅调整并行计算节点的规模（world size），实现通信开销的有效隔离与精确观测。

\textbf{基准组：}配置单机单卡训练环境，world size取1。该配置不包含跨设备通信开销，其性能数据作为纯计算性能的参照基准，用于量化分布式训练引入的额外通信成本。

\textbf{实验组：}配置单机双卡训练环境，world size取2，并启用FSDP全分片策略。该配置包含完整的梯度归约Reduce-Scatter与参数广播All-Gather通信原语，用于模拟分布式训练中的典型通信行为。

实验环境与关键参数配置如表~\ref{tab:bottleneck-config}~所示。

\begin{table}[htbp]
\caption{通信瓶颈分析实验参数配置}
\label{tab:bottleneck-config}
\vspace{0.5em}\centering\wuhao
\begin{tabular}{lll}
\toprule
参数项 & 配置值 & 说明 \\
\midrule
实验环境 & Ubuntu 24.04, 2$\times$A100 (40GB) & 操作系统与GPU硬件配置 \\
模型架构 & GPT-2 (Random Init) & 采用随机初始化以聚焦通信与计算特性 \\
训练精度 & FP32 & 全精度训练，避免混合精度对通信的干扰 \\
全局批大小 & 8 (Per-Device) & 确保单设备负载一致 \\
序列长度 & 1024 & 模拟典型长序列训练场景 \\
Profiler调度 & Wait=1, Warmup=1, Active=50 & 确保采集到稳定的训练迭代周期 \\
训练数据集 & Wikipedia En & 使用Wikipedia词条作为训练语料 \\
\bottomrule
\end{tabular}
\end{table}

\section{通信开销的定量分析}

基于Profiler采集的Trace数据，本研究对单个训练步内的计算流与通信流进行了细粒度的离线统计与分析。

\subsection{单卡与双卡训练耗时对比}

数据分析表明，双卡训练的平均Step耗时显著高于单卡基准，且该增量主要源于引入的通信开销。在单卡训练配置下，各模型规模的通信时间占比均处于极低水平，表明单步耗时主要由计算操作与本地I/O开销主导。而在双卡训练场景中，跨卡通信开销随模型规模增大呈现显著放大趋势：

（1）GPT2-small模型中，双卡单步耗时较单卡增加约50\%，其中通信占比达50.5\%。

（2）GPT2-medium模型中，双卡耗时增量进一步扩大，通信占比提升至52.0\%。

（3）GPT2-large模型中，双卡单步耗时较单卡增加超过100\%，通信占比高达61.2\%。

这一趋势清晰表明，随着模型参数规模增长，双卡训练中跨卡通信的相对开销显著提升，成为制约模型分布式训练性能的关键因素。单卡训练的单步耗时随模型规模近似线性增长，反映了计算成本在单设备训练中的主导地位。而双卡训练相对于单卡的额外开销几乎完全来源于跨卡通信，且通信时间占比随模型规模增大而持续升高。

\subsection{通信算子分解与重叠分析}

为更加准确地定位FSDP训练框架下的通信热点，本文进一步对双卡场景下的通信算子进行了分解分析。以GPT2-medium的训练数据为例：

在通信算子耗时构成方面，All-Gather操作占通信总耗时的61.7\%，Reduce-Scatter操作占38.3\%。在通信与计算重叠特性方面，暴露通信（Exposed/Blocking，即阻塞计算的通信）占比42.6\%，而可与计算重叠（Overlapped/Hidden）的通信占比57.4\%。

尽管All-Gather的绝对耗时更长，但结合FSDP的反向传播机制与通信暴露比例分析可发现：Reduce-Scatter通常位于梯度计算完成后梯度通信的关键节点上，其执行直接阻塞后续计算步骤，阻塞效应更为显著；而All-Gather主要发生在参数获取阶段，其通信可以通过预取与计算重叠机制被掩盖。

此外，从优化可行性角度看，All-Gather传输的是模型权重，对其进行有损压缩往往会引入显著的精度风险；而Reduce-Scatter传输的是梯度张量，其压缩空间更大且精度影响可控。因此，本研究将后向传播之后的梯度张量传输环节确立为通信压缩优化的首要对象，以在保证模型精度的前提下最大化性能收益。

\section{本章小结}

本章通过构建单卡与双卡GPT-2模型训练对比实验，利用PyTorch Profiler性能分析工具对FSDP分布式训练过程中的通信开销进行了精确量化与特征分析。实验结果表明：（1）通信开销随模型规模增大而显著增长，在GPT2-large模型中通信占比已达61.2\%；（2）Reduce-Scatter操作位于反向传播关键路径上，是通信压缩优化的首要目标；（3）All-Gather通信可通过预取机制部分掩盖，而Reduce-Scatter的阻塞效应更为显著。这些发现为后续通信压缩算法的设计与优化提供了坚实的实证依据。

